\documentclass[11pt,a4paper]{article}
\usepackage[T1]{fontenc}
\usepackage{lmodern,microtype}
\usepackage[margin=25mm]{geometry}
\usepackage{amsmath,amssymb,amsthm,mathtools,booktabs,array,xcolor}
\usepackage[colorlinks=true,linkcolor=blue!45!black,citecolor=blue!45!black,urlcolor=blue!45!black]{hyperref}
\hypersetup{pdftitle={A common zero-free strip for Hecke L-functions}}
\setlength{\emergencystretch}{2em}
\numberwithin{equation}{section}
\newtheorem{theorem}{Theorem}[section]
\newtheorem{proposition}[theorem]{Proposition}
\newtheorem{lemma}[theorem]{Lemma}
\newtheorem{hypothesis}[theorem]{Hypothesis}
\theoremstyle{remark}
\newtheorem{remark}[theorem]{Remark}
\newcommand{\eps}{\varepsilon}
\newcommand{\K}{\mathcal K}
\newcommand{\R}{\mathcal R}
\newcommand{\B}{\mathcal B}
\newcommand{\N}{\mathrm N}
\newcommand{\one}{\mathbf 1}
\title{A common zero-free strip for Hecke $L$-functions}
\author{AI-assisted research draft\thanks{AI assistance: OpenAI ChatGPT was used for mathematical exploration and drafting; OpenAI Codex was used for editing and compilation.}}
\date{8 October 2026}

\hypersetup{pdfsubject={Research draft}}
\begin{document}
\maketitle
\begin{abstract}
We prove a zero-free strip of width $25/258$ for all finite-order
Hecke $L$-functions, conditional on the cited theta and reflection
inputs and the fixed-order large sieve.  The proof first treats Dedekind zeta
functions.  Each reflected row then pairs a character with its
conjugate, and one symmetric zero-free rectangle controls both
$L$-factors.  The required density estimate retains
the number of selected rows in a fourth $L$-moment and a second
mollifier moment.  Its exponent is determined by a quadratic
equation, and a quartic inequality gives the stated rational width.
The same argument reaches an exactly defined endpoint
$0.096906264531643757\ldots$.  The width is common to all fixed
fields; uniform implied constants in a varying field are not asserted.
\end{abstract}

\section{The theorem and the reflection input}

We take as analytic inputs the theta realization, reflected
expansion, and their estimates in~\cite[Sections~3--8 and
Appendix~A]{OAI}.  The construction there is presented over
cyclotomic fields.  We first verify the fixed-field generality
needed here and record the precise consequences used in the proof.
The large-sieve input is stated separately in Section~\ref{sec-moments}.
These cited analytic results are retained as inputs; their full
proofs are not reproduced or independently certified in this article.

\begin{theorem}\label{thm-main}
Under the theta and reflection inputs in Proposition~\ref{prop-theta}
and the fixed-order large sieve \eqref{eq-ls}, for every number
field $K$ and every finite-order Hecke character $\eta$ of $K$,
\[
 L_K(s,\eta)\ne0\qquad\left(\Re s>\frac{233}{258}\right),
\]
with the usual pole at $s=1$ permitted for the principal character.
More precisely, the conclusion holds in $\Re s>1-\Delta_\#$,
where $\Delta_\#>25/258$ is defined in
\eqref{eq-endpoint}; its numerical value is
$0.096906264531643757\ldots$.
\end{theorem}

All implied constants may depend on the fixed field, the local
data, and fixed auxiliary parameters.  No zero-free hypothesis is
part of the inputs.  The new estimates below are deduced from
the displayed identities and the large sieve.

\begin{proposition}[Fixed-field theta reflection]\label{prop-theta}
Under the cited theta and reflection inputs, the construction
of~\cite{OAI} applies to every fixed number field $F\supset\mu_{12}$
after fixing a finite set $S$ and its local data.  It supplies
sixth-power-free rows $u$, characters
$\chi_u$ of order dividing six, and a probe $I(X,Y,Z)$ with the
following properties.

The row norm is $q_u=\N u$, the actual primitive conductor satisfies
$Q_u\ll_F q_u$, and the principal inducing row is $u=1$.
After a fixed finite ray-class partition, the characters form the
residue-character family in~\cite{DF}.  For $X=Y=Z^{1-h}$,
$1/2<h<3/5$, the reflected row square and the additive square
mass satisfy
\begin{equation}\label{eq-reflection-input}
 \sum_m|B_m(Z)|^2\ll Z^\eps(R_0+R_0^2/Z),\quad R_0=XY,
 \qquad \sum_m|A_m|^2\ll X^{-1}+Y^{-1},
\end{equation}
and $|I|\le(\sum|A_m|^2\sum|B_m|^2)^{1/2}$.

For the principal target, the Poisson expansion is a sum of triple
Mellin integrals with scalar factor
\begin{equation}\label{eq-scalar}
 X^{1/2-z}Z^{\xi+z-1}Y^{w-1}q_u^{-z}
\end{equation}
and row Dirichlet series
\begin{equation}\label{eq-euler}
 \zeta_F^S(6z)\,
 \frac{L_F^S(w,\chi_u)}{L_F^S(\xi,\overline{\chi_u})}
 H_u(\xi,w,z).
\end{equation}
The local factors are those in~\cite[(7.16)--(7.20)]{OAI}.
The fixed Mellin weights in $w,z$ are rapidly decreasing on
vertical lines, with the Gaussian factor $e^{(\xi+z-1)^2}$.
The double principal
residue at $w=1$, $z=1/6$
is $c_0f_h(Z)$, where $c_0\ne0$ and
\begin{equation}\label{eq-main-residue}
 f_h(Z)=\frac1{2\pi i}\int_{(2)}
 Z^{s-1+(1-h)/2+h/6}e^{(s-5/6)^2}
 \frac{H_1(s)}{\zeta_F^S(s)}\,ds,
 \qquad H_1(s)=H_1(s,1,1/6).
\end{equation}
The expansions and contour truncations have the convergence and
vertical decay asserted in~\cite[Sections~7--8]{OAI}.
\end{proposition}

\begin{proof}
Enlarge $S$ to contain the places above six, the trace-character
ramification, and primes generating the ideal class group.  The
ring of $S$-integers is then principal.  The Poisson lattice in
\cite[Lemma~3.1]{OAI} uses adelic duality, the product formula,
and strong approximation over $F$.  Balanced generators follow
from compactness of the norm-one idele class group.
Write $[F:\mathbb Q]=2r$; all its archimedean places are complex.
For $\ell=2,3$, the local and global $S$-unit power-class spaces
in \cite[Lemma~3.2]{OAI} have dimensions
\[
 \dim_{\mathbb F_\ell}(F_S^\times/F_S^{\times\ell})
 =2|S_f|+2r,\qquad
 \dim_{\mathbb F_\ell}(U_S/U_S^\ell)=|S_f|+r.
\]
The trivial $S$-class group gives injectivity of the global image;
Hilbert reciprocity makes it isotropic.  Since $-1$ is a sixth
power in $F$, the pairing is alternating.  Thus the image is
Lagrangian and has the isotropic complement required by the
construction.  Its Kummer characters and conductor bounds use
only $\mu_6\subset F$ and class field theory.

For the cubic Gauss correction, put $E_0=\mathbb Q(\sqrt{-3})$
and $\Lambda=\Lambda_{E_0}\circ N_{F/E_0}$.  If
$\mathfrak p\mid\mathfrak q$ has residue degree $f$, Jacobi-sum
lifting gives $J_{\kappa(\mathfrak p)}=(-1)^{f-1}
J_{\kappa(\mathfrak q)}^f=-\pi^+(\mathfrak q)^f$.
This proves the identity $\gamma_2(\mathfrak p)^3=-\Lambda(\mathfrak p)$
used in \cite[Lemma~3.4]{OAI}.  At every complex place the relative
norm contributes the coordinate or its conjugate, so the angular
exponent is $+1$ or $-1$, with no radial part.  This supplies the
constant-term cancellation even when $F$ is not CM.
The quadratic correction uses the Weil-index product formula
and the same isotropic complement.

The metaplectic framework of~\cite[Section~0.2 and Chapter~II]{KP}
is over any global field containing the required roots of unity.
For the specialization $(\text{rank},\text{cover degree},c)=(2,3,2)$
in~\cite[Proposition~5.1 and Appendix~A]{OAI}, the scalar center,
exceptional quotient, spherical vector, and local Whittaker
conditions depend on the local fields and these parameters.
The adelic splitting uses Hilbert reciprocity.  Compatibility of
the global datum with rational lifts uses that an everywhere-local
cube is a global cube, which follows from Kummer theory and
Chebotarev.  The cited nonzero theta residue and Whittaker
functional have the same general-field scope.

All archimedean profiles are therefore the complex profiles used
in the source.  The ensuing reflection uses those profiles and
finite power-class partitions; its ideal counts and quadratic
large sieve hold over each fixed number field.  The separated
additive fractions use balanced generators and the product
formula, and the exterior local Gauss identities use
$\N\mathfrak p\equiv1\pmod{12}$.  No step here requires $F/\mathbb Q$
to be Galois or abelian.

To justify the changed probe scales, retain the dyads in the
proof of~\cite[Section~6.4]{OAI}, before its final use of a fixed
row-box exponent.  Let $P=\N\mathfrak h$ denote the powerful
part and $D=\N\mathfrak d$ the squarefree part.  The compact
row box gives $PD\ll R_0=XY$.  The local reflection and quadratic
large sieve give, with the source's row restrictions retained,
\[
 \sum_{\mathfrak h\sim P}\sum_{\mathfrak d\sim D}|B_{hd}(Z)|^2
 \ll Z^\eps P^{1/2}\left(D+\frac{D^2P}{Z}\right)
 \ll Z^\eps\left(\frac{R_0}{\sqrt P}
                    +\frac{R_0^2}{Z\sqrt P}\right).
\]
Summing dyads proves the first estimate in \eqref{eq-reflection-input}.
The retained column cutoff is $Z^\nu(8D^2P/Z)$; the decay order
is chosen after $\nu$ in~\cite[(6.40)--(6.45)]{OAI}.
If this cutoff is below one, the full partial-sum estimate (6.40)
is already a rapidly decaying tail, so no constant term is added.
The additive square norm in~\cite[Section~7.1]{OAI} is
$(XY+Y^2)/(XY^2)=X^{-1}+Y^{-1}$.
The Poisson identities and the two principal residues then give
\eqref{eq-scalar}--\eqref{eq-main-residue} with these scales.
All constants depend only on the fixed field and auxiliary data.
\end{proof}

\subsection{The local Euler domain}

The exact local sums give a domain larger than the sufficient
domain recorded in~\cite[(7.21)]{OAI}.  This enlargement is needed
for the lines used in the proof.

\begin{lemma}\label{lem-euler}
The correction $H_u$ is holomorphic and $\ll q_u^\eps$, uniformly
in the imaginary parts, on every closed subregion with positive
margins of
\begin{equation}\label{eq-domain}
\begin{gathered}
 \Re\xi>\tfrac12,\quad \Re w>0,\quad \Re z>0,\quad
 \Re(\xi+w)>1,\\
 \Re(w+6z)>1,\quad\Re(\xi+6z)>1,\quad
 \Re(\xi+w+6z)>2,\quad\Re(\xi+z)>\tfrac56.
\end{gathered}
\end{equation}
After a fixed enlargement of $S$, $H_1(s)$ is nonzero for
$\Re s\ge9/10$.
\end{lemma}

\begin{proof}
Write $p=\N\mathfrak p$, and use real parts for the exponents
in the following bounds.  At a prime not dividing $u$, the exact
factor in~\cite[(7.19)]{OAI} differs from one by at most a constant
times
\[
 p^{-w-6z}+p^{-\xi-6z}+p^{-\xi-w}
       +p^{1-\xi-w-6z}+p^{4-6\xi-6z}.
\]
Every term is $O(p^{-1-\delta})$ on a fixed inner region of
\eqref{eq-domain}.  At a prime with $j=v_{\mathfrak p}(u)$,
$1\le j\le5$, the exact boundary terms instead give
\begin{align*}
 |H_{u,\mathfrak p}-1|\ll{}&p^{4-6\xi-6z}+p^{1-\xi-w}
 +\one_{j=5}p^{3-6\xi}+\one_{j=2}p^{3/2-3\xi}\\
 &+\one_{j=3}(p^{2-3\xi-w}+p^{2-4\xi})
 +\one_{j=4}p^{5/2-4\xi-w}.
\end{align*}
To obtain this bound directly from~\cite[(7.17)]{OAI}, write the
theta exponent as $t=6a,6a+1,6a+3,6a+4$.  For the other modulus
exponents $0,1$, the boundary terms occur at $j+6m=t-1,t$,
respectively.  Away from these boundaries the character sums
vanish unless $t\equiv0,1\pmod6$, respectively.  The remaining
terms sum geometrically with ratio of absolute value
$p^{4-6\xi-6z}$; the displayed bound retains each boundary.
All its exponents are negative on \eqref{eq-domain}.
Thus the good-prime product converges normally and the bad-prime
product is bounded by
$\prod_{\mathfrak p\mid u}(1+Cp^{-\delta})\ll q_u^\eps$.
For $u=1$ every exterior prime is good.  Enlarge $S$ so that each
remaining correction has distance less than $1/2$ from one on
$\Re s\ge9/10$, $w=1$, $z=1/6$; the convergent logarithmic
product proves nonvanishing.
\end{proof}

\section{Moments on selected rows}\label{sec-moments}

Fix $F$ and its local data.  We suppress the superscript $S$ in
moment estimates and omit the principal inducing row.  The input
from~\cite[Theorem~1.1]{DF}, with character order six, is
\begin{equation}\label{eq-ls}
 \sum_{q_u\le U}\left|\sum_{\N\mathfrak n\le V}
 a_{\mathfrak n}\chi_u(\mathfrak n)\right|^2
 \ll(UV)^\eps\K(U,V)\sum|a_{\mathfrak n}|^2,
 \quad
 \K(U,V)=U+V+U^{5/6}V^{1/3}+U^{1/3}V^{5/6},
\end{equation}
initially for sixth-power-free coefficient ideals.  Fixed twists,
ray-class partitions and row subsets are permitted.

We will also use \eqref{eq-ls} for divisor-bounded coefficients
on arbitrary ideals.  First restrict the complete ideal to
$\N\mathfrak n\asymp V$ and write
$\mathfrak n=\mathfrak b\mathfrak a^6$.  On a real line $s$, its
squared coefficient mass for fixed $\mathfrak a$ is
$O(V^{1-2s+\eps}/(\N\mathfrak a)^6)$.  Minkowski's inequality
and the convergent sum in $(\N\mathfrak a)^{-3+\eps}$ give
\begin{equation}\label{eq-dyad}
 \sum_{q_u\le U}|D_u(s+it;V)|^2
 \ll(UV)^\eps\K(U,V)V^{1-2s}.
\end{equation}
Here $\chi_u(\mathfrak a)^6$ merely restricts to rows coprime
to $\mathfrak a$.  Smooth weights can be separated on conductor
and column dyads by Mellin inversion.

The approximate functional equation and \eqref{eq-dyad} give
the three full-family bounds needed later:
\begin{equation}\label{eq-standard-moments}
\begin{gathered}
 \sum|L(1/2+it,\chi_u)|^2\ll U^{1+\eps}T_0^C,\qquad
 \sum|L(1/2+it,\chi_u)|^4\ll U^{7/6+\eps}T_0^C,\\
 \sum|L(1/4+it,\chi_u)|^8\ll U^{3+\eps}T_0^C,
\end{gathered}
\end{equation}
where $T_0=2+|t|$ and $C$ is fixed.  Indeed, apply the approximate functional
equation to $L^m$, with $m=1,2,4$, on actual conductor dyads
$Q\ll U$; the polynomial lengths are $Q^{m/2}T_0^C$.
Maximizing the four terms of \eqref{eq-dyad} at their two column
endpoints gives the displayed powers.  On the dual side retain
the squared multiplier $Q^{m(1-2s)}$ until the endpoints have
been evaluated.  Its lower endpoint is at most the displayed
bound; its upper endpoint has the same conductor exponent as
the original polynomial.  Conjugation inside an absolute square
restores the residue-character family, and the root number is
a row phase of modulus one.

For $A\ge1$, put
\[
 M_A(s,\chi_u)=\sum_{\N\mathfrak n\le A}
       \mu_F(\mathfrak n)\chi_u(\mathfrak n)(\N\mathfrak n)^{-s}.
\]
The next lemma preserves the size of an arbitrary selected set
in both factors of the zero detector.

\begin{lemma}\label{lem-sparse}
Let $\R$ be a set of $R\ge1$ rows of norm at most $U$, and set
$c=2-2\tau$, where $7/12\le\tau\le2/3$.  At a common height $t$,
\begin{align}
 \sum_{u\in\R}|L(\tau+it,\chi_u)|^4
 &\ll U^\eps T_0^C
       \left(U^cR^{1-c}+U^{c+1/6}\right),\label{eq-sparse-fourth}\\
 \sum_{u\in\R}|M_A(\tau+it,\chi_u)|^2
 &\ll(UA)^\eps\left(U^cR^{1-c}+A^c+
 U^{5/6}A^{c-2/3}+U^{1/3}A^{c-1/6}\right),\label{eq-sparse-second}
\end{align}
where \eqref{eq-sparse-second} requires $H=U/R\le A$.
Fixed multiplicative constants or subpower factors in $H$ are harmless.
\end{lemma}

\begin{proof}
The part of a divisor-bounded polynomial of norm at most $H$
is $O(H^{1-\tau+\eps})$ in every row.  Its squared sum on
$\R$ is $RH^c=U^cR^{1-c}$, up to allowed losses.
For a column dyad $V>H$, \eqref{eq-dyad} is
$\K(U,V)V^{c-1}$.  Its row term decreases with $V$, and its
other slopes $c,c-2/3,c-1/6$ are nonnegative.  Evaluating at
$H$ and $A$ proves \eqref{eq-sparse-second}.

For \eqref{eq-sparse-fourth}, apply the approximate functional
equation to $L^2$ and first partition the actual conductor into
dyads $Q$.  Ignore polynomial height factors temporarily; both
polynomials have length $Q$.  For the main polynomial, $Q\le H$
gives the rowwise estimate $RQ^c\le U^cR^{1-c}$.  If $Q>H$,
split at $H$ as above.  Its other endpoint terms are at most
$U^c$, $U^{5/6}U^{c-2/3}$ and $U^{1/3}U^{c-1/6}$,
all bounded by $U^{c+1/6}$.

The dual multiplier has absolute size $Q^{c-1}$ before squaring,
and the dual coefficient line is $1-\tau=c/2$.  All column
slopes on that line are positive.  At its upper endpoint, after
the squared multiplier is inserted, the four large-sieve terms are
\begin{equation}\label{eq-actual-conductor}
 UQ^{c-1},\qquad Q^c,\qquad
 U^{5/6}Q^{c-2/3},\qquad U^{1/3}Q^{c-1/6}.
\end{equation}
If $Q\ge H$, the first is at most $UH^{c-1}=U^cR^{1-c}$
because its conductor exponent is negative; the others are at
most $U^{c+1/6}$.  If $Q<H$, the entire dual polynomial,
including its multiplier, has squared rowwise bound $RQ^c$.
This is again at most $U^cR^{1-c}$.  Thus no negative conductor
power has been replaced prematurely by a power of $U$.
The smooth weights, sixth-power removal and dual conjugation are
handled as in \eqref{eq-dyad}; they cost only the displayed
polynomial height and subpower factors.
\end{proof}

\section{The density estimate}\label{sec-density}

For $T\ge1$, let $R(\sigma,T;U)$ count rows whose inducing nonprincipal
$L$-function has a zero in $\Re s\ge\sigma$, $|\Im s|\le T$.
We first state the detector in a form that retains the selected set.

\begin{lemma}\label{lem-detector}
Choose $1\le A<B$, $1/2\le\tau<\sigma<1$, with
$\sigma-\tau$ bounded below.  Apart from factors
$(UABT)^\eps T^C$, the polynomial-tail alternative for the
selected rows is bounded by
\begin{equation}\label{eq-tail}
 \B_\sigma=B^{2-2\sigma}+UA^{1-2\sigma}
 +\sup_{A\le V\le B}V^{1-2\sigma}
       (U^{1/3}V^{5/6}+U^{5/6}V^{1/3}).
\end{equation}
If $R$ selected rows satisfy the integral alternative, then
\begin{equation}\label{eq-detector}
 R^3\ll(UABT)^\eps T^C B^{4(\tau-\sigma)}
 \sup_{|v|\le T+C\log(UABT+3)}
 \left(\sum_{\R}|L(\tau+iv)|^4\right)
 \left(\sum_{\R}|M_A(\tau+iv)|^2\right)^2.
\end{equation}
\end{lemma}

\begin{proof}
Select one zero $\varrho_u=\beta_u+i\gamma_u$ in each row.
Mellin inversion and a shift to $\Re z=\tau-\beta_u$ give
\begin{align*}
 e^{-1/B}+\sum_{\N\mathfrak n>A}
 a_A(\mathfrak n)\chi_u(\mathfrak n)
 (\N\mathfrak n)^{-\varrho_u}e^{-\N\mathfrak n/B}
 =\frac1{2\pi i}\int_{(\tau-\beta_u)}
 L(\varrho_u+z)M_A(\varrho_u+z)\Gamma(z)B^z\,dz,
\end{align*}
where $a_A(\mathfrak n)=\sum_{\mathfrak d\mid\mathfrak n,
\,\N\mathfrak d\le A}\mu_F(\mathfrak d)$.
The crossed Gamma residue vanishes at $\varrho_u$; the omitted
principal row is the only possible $L$-pole.  Each row therefore
has a large integral or a large tail.  Gamma decay truncates the
integral at logarithmic distance from $\gamma_u$.  Substituting
$v=\gamma_u+\Im z$ puts all rows on a common height integral,
with a uniform envelope because $\beta_u-\tau$ stays away from
zero and one.  Summing and applying H\"older with exponents
$4,2,4$ gives
\[
 R\ll (UABT)^\eps T^C B^{\tau-\sigma}R^{1/4}
 \sup_v\left(\sum_{\R}|L|^4\right)^{1/4}
             \left(\sum_{\R}|M_A|^2\right)^{1/2},
\]
which proves \eqref{eq-detector}.

For the tail, dyadic subdivision and partial summation remove
the selected $\beta_u$.  The one-dimensional Sobolev inequality
replaces each row-dependent height supremum by an integral of
$|D_u|^2+|D'_u|^2$ on a slightly larger interval.  Binary blocks
handle the partial sums; derivatives and blocks cost logarithms.
At a common ordinate apply \eqref{eq-dyad}.  Maximizing its
two pure terms at $A,B$ gives \eqref{eq-tail}; the mixed terms
are left as a supremum.  This is the detector argument
of~\cite[Section~5]{BGL}, with the kernel \eqref{eq-ls}.
The exponentially decaying part beyond $B$ changes only allowed
losses.  At least half the rows satisfy one of the alternatives.
\end{proof}

\begin{proposition}\label{prop-density}
For $1/6\le d\le1/3$, let $x=x(d)$ be the unique positive root of
\begin{equation}\label{eq-quadratic}
 \Phi_d(x)=6d(6d^2-15d+7)x^2+(54d^2-39d+4)x+2d-2=0.
\end{equation}
Then, for some fixed $C$,
\begin{equation}\label{eq-density}
 R(2/3+d/2,T;U)\ll U^{k(d)+\eps}T^C,
 \qquad k(d)=\frac56-dx(d).
\end{equation}
\end{proposition}

\begin{proof}
The quadratic coefficient is positive and the constant is
negative.  The evaluations
\[
 \Phi_d(5/7)=\frac{(6d-1)(150d^2-35d-42)}{49}\le0,
 \qquad \Phi_d(1)=1+(1-3d)(1+8d-12d^2)\ge1
\]
give $5/7\le x<1$.  Put
\begin{equation}\label{eq-parameters}
 y=\frac{1-2dx}{1-2d},\qquad
 D_0=x(1-d)-\frac16,\qquad
 c=\frac{x(2/3-d)}{D_0},\qquad \tau=1-\frac c2.
\end{equation}
Here $D_0\ge13/42$, and
\[
 c-\frac23=\frac{1-3dx}{9D_0}\ge0,
 \qquad \frac56-c=\frac{6x(1+d)-5}{36D_0}\ge0.
\]
Also $1\le y\le11/7$, $x\le y$, and
$\sigma-\tau\ge1/12$ for $\sigma=2/3+d/2$.

Take $A=U^x$, $B=U^y$ and write $q=2\sigma-1=d+1/3$.
The two mixed terms in \eqref{eq-tail} have exponents
$1/3+(5/6-q)y$ and $5/6+(1/3-q)x$, both equal to $k$.
The differences from the other two exponents $1-qx$ and
$(1-q)y$ are $(2x-1)/6$ and $(2-y)/6$, both nonnegative.

At $R=U^k$, the condition $H\le A$ follows from $D_0>0$.
The row term in \eqref{eq-sparse-second} balances its first
mixed term because
\begin{equation}\label{eq-row-balance}
 c+(1-c)k=\frac56+(c-2/3)x.
\end{equation}
Its gaps from the other two terms are $5/6-2x/3$ and
$(1-x)/2$.  The row term in \eqref{eq-sparse-fourth} is no
larger than $U^{c+1/6}$: indeed
\[
 D_0-(2x-1)k=2(1-x)(1/3-dx)\ge0
 \quad\Longrightarrow\quad (1-c)k\le1/6.
\]
The exponent on the right of \eqref{eq-detector} is therefore
$4(\tau-\sigma)y+(c+1/6)+2\{c+(1-c)k\}$.
It equals $3k$, since
\begin{equation}\label{eq-final-balance}
 (1+2c)k=3c+\frac16+(4/3-2c-2d)y
\end{equation}
reduces exactly to \eqref{eq-quadratic} after \eqref{eq-parameters}
is substituted.

This verification bounds the unknown $R$ self-consistently.
For $R\ge U^k$, the condition $H\le A$ remains true.  Relative
to their values at $U^k$, each sparse moment grows by at most
$(R/U^k)^{1-c}$.  Thus \eqref{eq-detector} implies
\[
 (R/U^k)^3\ll U^\eps T^C(R/U^k)^{3(1-c)}.
\]
Since $c\ge2/3$, this gives \eqref{eq-density}.  For $R<U^k$
the claim is immediate.  The tail alternative was already bounded.
\end{proof}

\subsection{The exterior density ranges}

For $0\le d\le1/2$ set
\begin{equation}\label{eq-exterior}
 D=2+9d-6d^2,\quad
 \widehat x=\frac{8d+1}{D},\quad
 \widehat y=\frac{11d+2}{D},\quad
 \widehat k=\frac56-d\widehat x.
\end{equation}
Then $R(2/3+d/2,T;U)\ll U^{\widehat k+\eps}T^C$.
To see this, use \eqref{eq-detector} on $\tau=1/2$ with
$A=U^{\widehat x}$, $B=U^{\widehat y}$.  Here
$1/2\le\widehat x\le1$, $1\le\widehat y\le3/2$.
The full fourth moment is $U^{7/6}$, and the second mollifier
moment is at most $U^{5/6+\widehat x/3}$.  After taking the
cube root, the integral exponent is
\[
 \frac{17}{18}+\frac{2\widehat x}{9}
             -\left(\frac29+\frac{2d}{3}\right)\widehat y.
\]
Substitution makes it and the two mixed tail exponents equal to
$\widehat k$; the other gaps are again
$(2\widehat x-1)/6$ and $(2-\widehat y)/6$.

For $1/2\le\sigma\le2/3$, the row count at the endpoint and
\cite[Corollary~1.5]{DF} give, with $q=2\sigma-1$,
\begin{equation}\label{eq-low-density}
 R(\sigma,T;U)\ll U^{1-q/(4-6q)+\eps}T^C.
\end{equation}
For $11/12\le\sigma<1$, we use
\begin{equation}\label{eq-high-density}
 R(\sigma,T;U)\ll
 U^{(1-\sigma)(84\sigma-53)/6+\eps}T^C.
\end{equation}
For completeness, apply H\"older with exponents $4,4,2$ in the
integral alternative at $\tau=1/2$.  The fourth mollifier moment
is bounded by $\K(U,A^2)$, so the resulting detector is
\[
 R\ll(UABT)^\eps T^C
 \left(B^{1-2\sigma}U^{7/12}\K(U,A^2)^{1/2}+\B_\sigma\right).
\]
Set $q=2\sigma-1$, $x=(7q-3)/2$, $y=x+7/12$,
$A=U^x$, $B=U^y$.  Since $q\ge5/6$, the kernel exponent is
$2x$.  The integral term, the $B$ term and the larger mixed
tail term have exponent $(1-q)(42q-11)/12$; the other terms
are smaller.  This proves \eqref{eq-high-density}.

\section{The density envelope and the zero bins}

Use $x(d)$ from \eqref{eq-quadratic} on $[1/6,1/3]$ and
$\widehat x(d)$ from \eqref{eq-exterior} on its complement
in $[0,1/2]$.  Denote this piecewise choice by $x_*(d)$ and put
$k=5/6-dx_*(d)$, $r(d)=d(1-x_*(d))$.

\begin{lemma}\label{lem-envelope}
For $2/3\le\alpha\le11/12$, with $d=2\alpha-4/3$,
\begin{equation}\label{eq-envelope}
 k+2\alpha-\frac12=\frac53+r(d)
 \le\frac{43}{25}-\frac{m_0}{101250},
 \qquad m_0=\frac{247341}{71872}>0.
\end{equation}
Furthermore, if
\begin{equation}\label{eq-endpoint}
 r_\#=\max_{1/6\le d\le1/3}d(1-x(d)),\qquad
 h_\#=\frac1{5/3+r_\#},\qquad \Delta_\#=\frac{h_\#}{6},
\end{equation}
then $k+2\alpha-1/2\le1/h_\#$ throughout the same interval and
$25/43<h_\#<57/98<7/12$.
\end{lemma}

\begin{proof}
For the central interval test $x_0=1-4/(75d)$.  Direct substitution
and $t=d-1/4$ give
\begin{align*}
 75^2d\Phi_d(x_0)
 &=202500d^4-224100d^3+66501d^2-3690d-528\\
 &=202500t^4-21600t^3-\frac{51273}{2}t^2
                              +198t-\frac{303}{64}\\
 &\le-\frac{30321}{4}t^2+198t-\frac{303}{64}
 \le\frac{39204}{30321}-\frac{303}{64}=-m_0,
\end{align*}
where $|t|\le1/4$ suffices for the first inequality.  Thus
$x_0<x(d)$.  Between these two values, $\partial_x\Phi_d<18$;
the mean value theorem gives
$d(x-x_0)\ge m_0/101250$ and hence \eqref{eq-envelope}.

For the exterior parameters, the exact identities
\[
 \frac1{20}-d(1-\widehat x)
 =\frac{(1-6d)(2+d-20d^2)}{20D}\ge0\quad(0\le d\le1/6),
\]
\[
 \frac2{39}-d(1-\widehat x)
 =\frac{(3d-1)(78d^2+9d-4)}{39D}\ge0\quad(1/3\le d\le1/2)
\]
give stronger bounds than \eqref{eq-envelope}.  Finally,
$\Phi_{1/4}(15/19)=81/5776>0$ implies $r_\#>1/19$.
This exceeds both exterior bounds, so the central maximum controls
the entire interval.  Combining $1/19<r_\#<4/75$ gives the
asserted range of $h_\#$.
\end{proof}

\subsection{The numerator and denominator}

Fix a small grid width $\lambda>0$, a dyad $q_u\asymp U$, and
a small fixed height exponent $\nu>0$.  Put a nonprincipal row
in the bin $[\alpha,\alpha+\lambda]$ when its rightmost zero in
$\Re s\ge1/2$, $|\Im s|\le4U^\nu$ has real part in that interval.
Rows with no such zero form the floor group.

For a binned row, on the fixed real-part range used below,
\begin{equation}\label{eq-both}
 |L_F^S(s,\chi_u)|+|L_F^S(s,\chi_u)|^{-1}\ll U^\eps
 \qquad(\Re s\ge\alpha+2\lambda,\ |\Im s|\le U^\nu).
\end{equation}
Indeed choose nested discs centered at $2+it$ whose largest disc
stays to the right of $\alpha+\lambda$ and within the zero-free
rectangle.  A branch of $\log L$ is bounded on a small disc in
absolute convergence.  The polynomial conductor-height bound and
Borel--Carath\'eodory bound it by $O(\log U)$ on a larger inner
disc.  Three-circles interpolation on an intermediate disc reaching
$\alpha+2\lambda+it$ gives $O((\log U)^\theta)$ with fixed
$\theta<1$.  Exponentiating gives both parts of \eqref{eq-both}.
Since $L_F^S(s,\overline{\chi_u})=\overline{L_F^S(\bar s,\chi_u)}$,
the symmetric rectangle gives the same bounds for $\overline{\chi_u}$,
in particular the reciprocal bound for the denominator in \eqref{eq-euler}.
The same argument applies to the floor with left boundary $1/2$.

For $\alpha\le11/12$ choose
\begin{equation}\label{eq-bin-lines}
 \xi=\alpha+4\lambda,\qquad
 w=1-\alpha-2\lambda,\qquad z=1/3+\lambda.
\end{equation}
The primitive functional equation sends the numerator at $w+it$
to the conjugate character at $1-w-it=\alpha+2\lambda-it$.
The fixed local factors at $S$ are harmless, so \eqref{eq-both} gives
\[
 |L_F^S(w+it,\chi_u)|
 \ll U^{\alpha-1/2+2\lambda+\eps}(2+|t|)^C.
\]
Here the conductor exponent is nonnegative, so $Q_u\ll U$ can
be used directly.  Choose $\nu$ small enough to absorb the height
factor.  A bin containing $O(U^{k+\eps})$ rows therefore has
absolute ratio sum bounded by
\begin{equation}\label{eq-ratio-sum}
 U^{\ell+O(\lambda)+\eps},\qquad
 \ell=k+\alpha-\tfrac12.
\end{equation}
The numerator and denominator heights can vary independently in
their central ranges.  Both are controlled by the same zero-free
rectangle; no independence of a selected set is required.

\section{Contour estimates and completion of the proof}

Fix $25/43\le h\le h_\#$ and write $a=1-h$, $X=Y=Z^a$.
Since $2a<1$, \eqref{eq-reflection-input} gives
\begin{equation}\label{eq-probe-bound}
 I(Z^a,Z^a,Z)\ll Z^{a/2+\eps}.
\end{equation}
For rows $U=Z^t$ whose absolute ratio sum has exponent $\ell$,
\eqref{eq-scalar} gives the relative exponent
\begin{equation}\label{eq-ledger}
 \mathcal E-a/2=\xi-1+hz+a(w-1)+t(\ell-z).
\end{equation}
Take a grid containing $1/2,2/3,3/4,5/6,11/12$ and sufficiently
small $\lambda$, and split at $U=Z^{1/6}$ and $U=Z^{h+\lambda}$.
All intermediate rows use $z=1/3+\lambda$.

The following table gives the complete bounds.  Lines and exponents
in the table omit $O(\lambda)$ offsets; the actual bin lines are
\eqref{eq-bin-lines}.  Every bound is relative to $Z^{a/2}$.
\begin{center}
\small
\begin{tabular}{@{}llll@{}}
\toprule
Range & $(\xi,w,z)$ & Ratio exponent $\ell$ & Relative exponent\\
\midrule
Principal $w$ remainder & $(1,2/3,1/6)$ & $0$ & $\le-1/24$\\
Principal $z$ remainder & $(1,1,0)$ & $0$ & $0$\\
$U\le Z^{1/6}$ & $(1,1/100,1/3)$ & $299/200$ & $\le-11/450$\\
Floor & $(1/2,1/2,1/3)$ & $1$ & $\le-1/8$\\
$1/2\le\alpha\le2/3$ & $(\alpha,1-\alpha,1/3)$ & $k+\alpha-1/2$ & $\le-1/36$\\
$2/3\le\alpha\le11/12$ & $(\alpha,1-\alpha,1/3)$ & $k+\alpha-1/2$ & $\le0$\\
$11/12\le\alpha<1$ & $(\alpha,1/4,1/3)$ & $7k/8+3/8$ & $\le-1/144$\\
$U>Z^{h+\lambda}$ & $(4,4,6/\lambda+1)$ & $1$ & $\le-2h$\\
\bottomrule
\end{tabular}
\end{center}

Here are the contour choices and verifications.  For the principal
row first take $\xi=1+\lambda$, move $w$ to $2/3$ with
$z=1/6+\lambda$, and take the pole at $w=1$.  Inside its residue
move $z$ to $\lambda$ and take the pole at $1/6$.  The first
remainder has relative exponent $h/2-1/3\le-1/24$, and the
second has exponent $O(\lambda)$.  These are sequential moves;
the pair $w=2/3$, $z=\lambda$ is never required.

For small rows take $\xi=1+\lambda$, $w=1/100$.  Pointwise
convexity and the row count give $\ell=299/200$.  Substitution
at $t=1/6$ gives the table, since $h\le7/12$.  For the floor,
take $\xi=1/2+2\lambda$, $w=1/2$ and use \eqref{eq-both}
and the first moment from \eqref{eq-standard-moments}.
Its exponent at $t=h$ is $-1+3h/2\le-1/8$.

For bins below $11/12$, \eqref{eq-ledger} and
\eqref{eq-ratio-sum} simplify at $t=h$ to
\begin{equation}\label{eq-crucial-exponent}
 -1+h\left(k+2\alpha-\frac12\right).
\end{equation}
In the low range put $q=2\alpha-1$.  The function
$1-q/(4-6q)+q+1/2$ increases on $[0,1/3]$ and has maximum
$5/3$, giving $-1/36$.  In the middle range
Lemma~\ref{lem-envelope} makes \eqref{eq-crucial-exponent}
nonpositive for $h\le h_\#$.  At $h=25/43$ it is at most
$-m_0/174150<-1/60000$.

For the highest bins take $\xi=\alpha+2\lambda$, $w=1/4$.
The reciprocal is controlled by \eqref{eq-both}; H\"older and
the eighth moment in \eqref{eq-standard-moments} give
$\ell=7k/8+3/8$, with $k$ from \eqref{eq-high-density}.
The relative exponent at $t=h$ is
$\alpha-7/4+h(9/8+7k/8)$.  Its derivative in $\alpha$ is
negative in this interval for $h\ge25/43$.  Its maximum is
$-5/6+17h/12\le-1/144$.

In every intermediate group $\ell-z>0$.  In the middle bins,
$\ell=1+d(1/2-x_*)\ge3/4$; in the highest bins
$\ell\ge3/8>1/3$.  The largest dyadic endpoint is consequently
$t=h+\lambda$, changing the bounds by $O(\lambda)$.
All contours satisfy \eqref{eq-domain}.  On the bin contours
\eqref{eq-bin-lines}, the potentially tight local Euler inequalities have margins
$\xi+w-1=2\lambda$ and $\xi+z-5/6\ge5\lambda$.

For large rows move to $\xi=w=4$ and $z=6/\lambda+1$.
Absolute convergence gives the row exponent one.  Since $z>1$,
the worst endpoint is $t=h+\lambda$; insertion in
\eqref{eq-ledger} gives $-2h$.

Finally, for each fixed grid choose $\nu$ sufficiently small
that the factors $T^C$ in the density and moment estimates,
with $T=4U^\nu$, cost only $U^{O(\lambda)}$.
Truncate the original Mellin contours before shifting their central
segments.  The horizontal joins lie in the zero-free rectangles
used for \eqref{eq-both}; the remaining tails stay on lines of
absolute convergence.  The Gaussian and rapidly decreasing
Mellin weights give any fixed power saving there.  Since
$U>Z^{1/6}$ for intermediate rows, these savings also control
the horizontal integrals in the $Z$ aspect.  The small-row and
large-row contours need only absolute convergence and convexity.
All fixed decay orders and local data are chosen before $Z$ grows.

The Poisson expansion, the residues and the table now yield
\[
 I(Z^a,Z^a,Z)=c_0f_h(Z)+O(Z^{a/2+C\lambda}),
 \qquad f_h(Z)\ll Z^{a/2+C\lambda}.
\]
For $0<Z\le1$, moving the line in \eqref{eq-main-residue}
arbitrarily far right gives arbitrary power decay.  Its Mellin
transform therefore continues
$e^{(s-5/6)^2}H_1(s)/\zeta_F^S(s)$ holomorphically to
$\Re s>1-h/6+C\lambda$.  By Lemma~\ref{lem-euler} its numerator
does not vanish in the required region.  Letting
$\lambda\downarrow0$ proves that $\zeta_F$ has no zeros in
$\Re s>1-h/6$, including $h=h_\#$; nonpositive limiting
exponents are sufficient for this open half-plane conclusion.

To finish, let $K$ be arbitrary.  The extension $K(\mu_{12})/K$
is finite abelian and its Dedekind zeta function factors into
$\zeta_K$ and nonprincipal, entire Hecke factors.  The strip
proved for the larger field therefore excludes zeros of
$\zeta_K$ away from one.  Next, any finite-order Hecke character
$\eta$ occurs in the zeta factorization of a finite abelian
extension $E/K$.  The strip for $\zeta_E$ excludes zeros of
$L_K(s,\eta)$ away from one.  Comparing the nonzero simple
residues of $\zeta_E$ and $\zeta_K$ gives nonvanishing of each
nonprincipal factor at one.  This proves Theorem~\ref{thm-main}.

\appendix
\section{The optimized endpoint}

Here $x(d)$ always denotes the central root in \eqref{eq-quadratic}.
Set $r=d(1-x)$ and
\begin{align*}
 Q(d,r)=d\Phi_d(1-r/d)
 ={}&36d^4-(36+72r)d^3+(5+126r+36r^2)d^2\\
 &+(2-45r-90r^2)d-4r+42r^2.
\end{align*}
Along the root, $Q_r=-\partial_x\Phi_d<0$.
On $1/6\le d\le1/3$, $0\le r\le4/75$,
\[
 Q_{dd}=432d^2-(216+432r)d+10+252r+72r^2
 \le-\frac{2622}{625}<0.
\]
Indeed this expression increases with $r$ and is convex in $d$;
its maximum is obtained by setting $r=4/75$ and testing the two
endpoints in $d$.  At every stationary point of $r(d)$,
$r''=Q_{dd}/\partial_x\Phi_d<0$, so there is at most one.
There is an interior maximum because
\[
 r(1/6)=1/21<1/19,\quad
 \Phi_{1/3}(16/19)=-28/361<0,\quad
 \Phi_{1/4}(15/19)=81/5776>0.
\]
Thus the maximum defining $r_\#$ is attained at a unique $d_\#$.

Its value can be enclosed using rational arithmetic alone.
For rational $d$, bisect the positive root $x(d)$ in $[5/7,1]$.
The sign of $r'(d)$ is the sign of
\[
 g(d,x)=(36d^3-42d)x^2
 +(72d^3-126d^2+84d-4)x+54d^2-37d+4.
\]
Rational interval evaluations of $g$ permit an outer bisection
of $d_\#$.  The bounds $\partial_x\Phi_d>3$ and
$|\partial_d\Phi_d|\le38$ give $|r'(d)|<6$ on this interval.
Hence a root interval $[d_-,d_+]$ and an interval evaluation of
$r$ at its midpoint enclose $r_\#$ after adding
$3(d_+-d_-)$ to the upper endpoint.  With 180 outer and 280
inner bisections, exact rational arithmetic gives
\begin{align*}
 0.253879034133118787&<d_\#<0.253879034133118788,\\
 0.053208384468374847&<r_\#<0.053208384468374848,\\
 0.096906264531643757&<\Delta_\#<0.096906264531643758.
\end{align*}
An exact-arithmetic implementation accompanies the article.
The numerical enclosure is not needed for the rational width
$25/258$ or for the exact variational definition \eqref{eq-endpoint}.

\begin{thebibliography}{9}
\bibitem{OAI}
OpenAI, \emph{Primitive roots for every admissible integer base},
4 October 2026, 92 pp., OpenAI Math, family 029.
\href{https://github.com/openai/math/blob/main/preprints/Primitive-roots-for-every-admissible-integer-base-October-4-2026/primitive-roots-all-integer-bases.pdf}{Official repository manuscript}.
\bibitem{DF}
A. de Faveri, \emph{Optimal large sieve for fixed order characters},
\href{https://arxiv.org/abs/2610.04045}{arXiv:2610.04045v1},
2 October 2026.
\bibitem{BGL}
V. Blomer, L. Goldmakher and B. Louvel,
\emph{$L$-functions with $n$-th-order twists},
Int. Math. Res. Not. \textbf{2014}, no.~7, 1925--1955.
\bibitem{KP}
D. A. Kazhdan and S. J. Patterson, \emph{Metaplectic forms},
Publ. Math. IH\'ES \textbf{59} (1984), 35--142.
\end{thebibliography}
\end{document}
